\scanpage{569}
\mspage{478}
\deleted{Pour tout} Soit $\underline S(\pi)_0$ une sous-catégorie réduite
de $\underline S(\pi)$, telle que $\underline S(\pi)_0\hookrightarrow\underline S(\pi)$ soit
une équivalence, de sorte que $\underline S(\pi)_0$ est
petite. Pour tout $x\in\operatorname{Ob}\underline S(\pi)_0$, soit
$i_x\in I$ \deleted{\unsure{un des indices qui}}\inserted{tel que} $x\in\operatorname{Ob}\underline S_{i_x}$. \unsure{Soit donc|Soit}
$J$ l'ens de tous les $i_x$, qui est
une petite partie de $I$, et soit
$\underline S_J$ la sous-catégorie \inserted{str. pl.} de $M$ ayant
pour \deleted{\unsure{ens. d'}} ens. d'objets $\bigcup_{j\in J}\operatorname{Ob}\underline S_j$.

\deleted{\unsure{On a}} Donc $\underline S(\pi)_0$ est contenue
dans $\underline S_J$, donc aussi $\underline S(\pi)$ puisque
$\underline S_{\deleted{\unsure{i}}J}$ est strictement pleine :
\[
\underline S(\pi)\subset\underline S_J.
\]
D'ailleurs $\underline S_J$ est accessible dans $M$
(loc. cit), et est fermée dans $M$, comme
réunion de sous-catégories fermées, \uline{a fortiori}
\uline{elle satisfait à} $L_\pi$. Donc, comme $\underline S(\pi)\subset\underline S_J$,
le foncteur \inserted{(\unsure{ess.} surjectif)} $\underline{\operatorname{Ind}}_\pi(\underline S(\pi))\longrightarrow\underline S$ se
factorise par $\underline S_J$. Donc $\underline S_J\hookrightarrow\underline S$ est
ess. surjectif, et comme $\underline S_J$ est str.
pleine, on en conclut $\underline S_J=\underline S$, \unsure{qed|cqfd}.
% archive sheet number: 568
